% Lettre classique, sans questions ni sections, police Helvetica comme le CV.
\documentclass[12pt,a4paper]{article}
\usepackage[a4paper,top=1.9cm,bottom=1.9cm,left=2.4cm,right=2.4cm]{geometry}
\usepackage[utf8]{inputenc}
\usepackage[T1]{fontenc}
\usepackage[scaled=1.0]{helvet}
\renewcommand{\familydefault}{\sfdefault}
\usepackage[hidelinks]{hyperref}
\usepackage{xcolor}
\definecolor{muted}{HTML}{555F67}
\setlength{\parindent}{0pt}
\setlength{\parskip}{0.65em}
\linespread{1.25}
\emergencystretch=2em
\pagestyle{plain}
\clubpenalty=10000
\widowpenalty=10000
\brokenpenalty=10000
\begin{document}
\thispagestyle{empty}
\textbf{Stéphane KPOVIESSI}\\
{\small Cergy, Île-de-France\\
+33 7 48 59 47 13\\
\href{mailto:oastephanekpoviessi@gmail.com}{oastephanekpoviessi@gmail.com}}

\vspace{0.4em}
Comité d'admission\\
Faculté des sciences de l'administration\\
Université Laval

\vspace{0.3em}
\textbf{Objet : Candidature au MBA en analytique d'affaires -- Automne 2027}

Madame, Monsieur,

Ingénieur diplômé de JUNIA ISEN, spécialisé en Big Data et Data Science, je souhaite intégrer le MBA en analytique d'affaires de l'Université Laval. Mes expériences m'ont appris à produire des analyses fiables pour les équipes. Je souhaite désormais mieux comprendre comment une organisation définit ses priorités, évalue un investissement et mobilise ses ressources pour décider.

Mon cursus d'ingénieur, suivi de 2020 à 2025 au Bénin, au Maroc puis en France, m'a donné des bases quantitatives solides. Après deux années préparatoires scientifiques, je les ai approfondies en traitement des données, en visualisation et en apprentissage automatique à JUNIA ISEN. Ce parcours dans trois pays m'a aussi appris à m'adapter à des méthodes de travail et à des interlocuteurs différents.

Chez Allianz France, la migration d'un datamart d'assurance m'a montré que la fiabilité d'un indicateur dépend autant des règles métier que des traitements techniques. Chez Lyne \& Frères, les analyses commerciales et opérationnelles m'ont rapproché des préoccupations d'une direction : comprendre un écart, suivre l'activité et disposer d'informations utilisables. Chez Sylvagreg, la reconstruction 3D par apprentissage profond et vision par ordinateur a renforcé ma pratique de l'expérimentation en IA. Ces expériences m'ont donné envie d'apprendre à formuler des recommandations tenant compte des objectifs et des contraintes de l'entreprise.

Je choisis l'Université Laval pour l'articulation de son MBA entre formation générale en gestion et spécialisation en analytique d'affaires. Je souhaite développer une lecture financière et stratégique des problèmes : apprécier la valeur d'un projet, comparer des options et comprendre leurs conséquences sur la performance. La place accordée à l'aide à la décision et à une activité de synthèse pratique correspond à mon souhait de travailler sur des situations concrètes. Je m'intéresse particulièrement à la manière de relier les indicateurs commerciaux et opérationnels aux décisions d'allocation des ressources.

\begin{samepage}
À court terme, je vise des fonctions d'analyse de la performance, de conseil en transformation ou d'analytique au sein d'une direction métier. Je souhaite contribuer aux choix de développement et aux arbitrages financiers, en associant compréhension des données et vision de l'entreprise. À moyen terme, mon ambition est de prendre la responsabilité de projets réunissant des équipes techniques et de gestion, puis d'évoluer vers des responsabilités de pilotage stratégique.\par
\end{samepage}

Depuis septembre 2026, je développe Seshat, un projet de plateforme d'IA générative associant recherche documentaire et automatisation. Structurer ce projet demande de comprendre les besoins des entreprises et d'arbitrer des ressources. Je souhaite acquérir au MBA les méthodes nécessaires à son développement et, à long terme, à celui de Lyne \& Frères. Mes responsabilités en communication puis à la vice-présidence de JUNIA ALUMNI m'ont aussi initié à la coordination et au travail collectif ; je souhaite approfondir ces compétences de gestion et de leadership.

Je souhaite apporter aux échanges à Laval ma perspective d'ingénieur, mes expériences dans des organisations de tailles différentes et les questions concrètes soulevées par mon projet entrepreneurial. Je vous remercie de l'attention portée à ma candidature et me tiens à votre disposition pour tout renseignement complémentaire.

Veuillez agréer, Madame, Monsieur, l'expression de mes salutations distinguées.

\vspace{0.5em}
\textbf{Stéphane KPOVIESSI}
\end{document}
